% !TeX root = ../../Book5.tex
% 纲要标题：### 27.7 推前形式与解析指标的比较
% 纲要位置：几何学计划纲要.md，第 6186 行。
% 正文状态：未撰写；本文件目前是框架。

% 纲要内容开始
% 
% 先准确使用高阶直像有限局部自由的固有光滑版本，再说明更一般表述需要的对象；不对任意奇异态射直接套用向量丛公式。[R11]
% 
% * 高阶直像的交错和与普通直像的区别
% * 固有光滑态射上相对 Todd 类的公式
% * 高阶直像不局部自由时的 K 群或完美复形修正
% * 一般 Grothendieck–Riemann–Roch 的外部理论位置
% * 复数域下代数、解析及拓扑 Chern 特征的比较
% 
%   解析指标比较只在复数域及已具备相应比较定理时使用；任意基域的代数公式不能由复解析证明不加说明地转移。
% 
% * 与 II.24.4、III.17、IV.20 的指标和 GAGA 路线衔接
% 
% ---
% 
% 纲要内容结束

\section{推前形式与解析指标的比较}
\label{g5:sec:2707}

待撰写

\subsection{高阶直像的交错和与普通直像的区别}
\label{g5:subsec:2707:01}
待撰写

\subsection{固有光滑态射上相对 Todd 类的公式}
\label{g5:subsec:2707:02}
待撰写

\subsection{高阶直像不局部自由时的 K 群或完美复形修正}
\label{g5:subsec:2707:03}
待撰写

\subsection{一般 Grothendieck–Riemann–Roch 的外部理论位置}
\label{g5:subsec:2707:04}
待撰写

\subsection{复数域下代数、解析及拓扑 Chern 特征的比较}
\label{g5:subsec:2707:05}
待撰写

\subsection{与 II.24.4、III.17、IV.20 的指标和 GAGA 路线衔接}
\label{g5:subsec:2707:06}
待撰写
% 纲要条目排为编号小节标题，并进入目录与书签；说明段仍保留于纲要注释中。
